\section{Chapter 13 Lower-Bound Basic Ideas Spine}

\paragraph{Status.}
This note is a synchronized natural-language export for task
\texttt{TEXTBOOK-PART-IV-CHAPTER-13-BASIC-LOWER-BOUND-SPINE}.  The semantic
and deterministic Chapter 13 slice passed its local Lean and repository gates
at commit \texttt{47ef4fe42679cd53e208aed540b6aa2bdb253aa8}.  It does not
claim the unit-variance Gaussian minimax lower bound in Theorem~13.1, whose
proof the source defers to Chapter~15.

\paragraph{Lean declarations.}
The compiled public surface consists of
\texttt{worstCaseExpectedRegret}, \texttt{minimaxExpectedRegret}, their three
order lemmas, \texttt{exists\_alternative\_le\_average},
\texttt{alternativeExpectedPullBudget\_le},
\texttt{exists\_leastExploredAlternative}, the two regret expressions, and
\texttt{max\_base\_changed\_regretLowerBound\_ge\_half\_sub\_error} with its
zero-error corollary \texttt{max\_base\_changed\_regretLowerBound\_ge\_half}.

\paragraph{Proof map.}
Let $R(\pi,\nu)$ take values in $[0,\infty]$.  The worst-case value of a policy
is the supremum over the explicit environment subtype, and the minimax value
is the infimum of these worst-case values over the explicit policy subtype.
The order lemmas follow from the corresponding complete-lattice introduction
and elimination rules.

For $m+1$ arms, split the finite sum at arm zero.  Nonnegativity of its
expected pull count and the exact total-pull identity imply that the sum over
the $m$ alternative arms is at most the horizon.  Finite averaging then gives
an $i\in\mathrm{Fin}(m)$ whose successor arm has expected count at most $n/m$.

For the two-environment algebra, write
\[
  B=\Delta(n-x),\qquad C=\Delta y.
\]
If $\Delta\geq 0$ and $x-y\leq e$, then
$B+C\geq\Delta(n-e)$.  Since both $B$ and $C$ are at most
$\max\{B,C\}$, it follows that
\[
  \frac{\Delta(n-e)}{2}\leq\max\{B,C\}.
\]
The half-horizon statement is the special case $e=0$ under $x\leq y$.
The transport comparison is an explicit premise: no likelihood-ratio, KL,
absolute-continuity, policy-consistency, or Gaussian result is inferred from
this deterministic theorem.

\paragraph{Verification boundary.}
The focused module, typed root-import canary, full Tests target, axiom print,
repository check, lean-verified website build, site checker, and desktop/mobile
browser inspection passed locally.  Theorem~13.1 remains planned until the
Chapter~14 history-law bridge and Chapter~15 minimax construction compile.
